\article{模 $p$ 的 Shimura 簇上的点}{Points on Shimura Varieties mod $p$}{J. S. Milne}{2}{165--184}

给定 $\Q$ 上带有一定附加结构的约化群（reductive group）$G$，可以构造相应的
复 Shimura 簇（Shimura variety）$S_\C$。多数情形已知它有定义在某个指定数域
$E$ 上的典范模型（canonical model）$S_E$。对 $E$ 的几乎所有有限素位 $v$，可将
$S_E$ 模该素位约化，得到有限域 $\F_{q_v}$ 上的非奇异簇 $S_v$。如 [3] 所述，为确定
$S_E$ 及其上局部常值层（locally constant sheaf）的 Hasse--Weil $\zeta$ 函数，须描述
$(S_v(\overline\F_{q_v}),\operatorname{Frob})$：前者为 $S_v$ 在
$\F_{q_v}$ 的代数闭包中的点集，后者为 Frobenius 映射，将坐标为
$(a_1,\ldots,a_m)$ 的点送到 $(a_1^{q_v},\ldots,a_m^{q_v})$。
所需描述应直接用群 $G$ 的数据来表达。

最近，Langlands [13] 对任意 Shimura 簇在足够好的素位处提出了这种描述的猜想。
在 [12] 中，他对能实现为某些模问题的粗模概形（coarse moduli schemes）的 Shimura 簇，给出了相当详细的证明纲要；这些模问题仅涉及
阿贝尔簇（abelian varieties）、弱极化（weak polarizations）、自同态（endomorphisms）及有限阶点。
因此 $G(\Q)$ 形如 $\Aut_B(H_1(A,\Q),\psi)$，其中 $B$ 是半单
$\Q$ 代数，包含一个作用于 阿贝尔簇 $A$ 的序（order）；$\psi$ 是 $A$ 的 Riemann 形式（Riemann form），其在
$\End(A)\otimes\Q$ 上的 Rosati 对合（Rosati involution）稳定 $B$。
$\Aut_B$ 指 $H_1(A,\Q)$ 的 $B$ 线性自同构 $g$，满足
$\psi(gu,gv)=\psi(u,\mu(g)v)$，其中 $\mu(g)$ 属于 $B$ 的中心内某个固定的、被
Rosati 对合固定的代数 $F$。此外还有一个 Hasse 原理（Hasse principle）假设。

更早 Ihara [8，猜想 1] 已在 $S$ 为 Shimura 曲线时提出类似猜想，并在
$G=\GL_2$ 时证明 [9，第 5 章]。当 $G=B^\times$，$B$ 为 $\Q$ 上四元数除代数（quaternion division algebra）时，
Morita [15] 对 $E=\Q$ 中所有不整除 $B$ 的判别式的素数 $p$ 证明了 Ihara 猜想。
他与 Shimura 对更一般的四元数代数都得到过部分结果（未发表）。最近 Ihara 对所有
Shimura 曲线及足够好的素位证明了其猜想（见 [11] 的公告）。Ihara 的证明基于
Eichler--Shimura 同余关系（congruence relations），而按 [10] 描述的 Morita 方法，似乎与
Langlands 的方法颇相似。

为介绍 Langlands 证明中的方法，本文讨论 $G$ 为全实数域上全不定四元数代数的乘法群这一情形。
第 1 节说明 $S_\C$ 自然参数化一族带附加结构的 阿贝尔簇。下一节介绍如何用
Artin 可表性判据（representability criteria）证明：存在 $\Q$ 上的簇
$S_\Q$，它是 $S_\C$ 的典范模型，且其模 $p$ 约化参数化特征 $p$ 下带附加结构的
阿贝尔簇族。由此，描述 $(S_p(\overline\F_p),\operatorname{Frob})$ 的问题便转化为研究这类带结构的阿贝尔簇。
第 5 节用有限域上 阿贝尔簇同源类（isogeny classes）的 Tate--Honda 分类确定该族中的同源类，第 6 节描述各个同源类。
由于需要 $p$ 可除群（$p$-divisible groups）及其 Dieudonné 模（Dieudonné modules），第 3 节先回顾这些内容。
\footnotetext{AMS（MOS）主题分类（1970）：主要 14K10；次要 14G10、14L05。原文版权：1979 年，美国数学会。}

\paragraph*{记号（Notation）。}
$F$ 为 $\Q$ 上次数 $d$ 的全实数域，$B$ 为 $F$ 上在所有无穷位处分裂的四元数除代数，
$b\mapsto b^*$ 为 $B$ 的正 $F$ 对合（positive $F$-involution），$O_B$ 为 $B$ 的极大序（maximal order）。

$G$ 为 $\Z$ 上群概形（group scheme），使对任意环 $R$，
$G(R)=(O_B^{\mathrm{opp}}\otimes R)^\times$，其中 $O_B^{\mathrm{opp}}$ 为 $O_B$ 的反代数（opposite algebra）。

$\A$ 为 $\Q$ 的阿代尔环（adele ring）；
\[
\A=\R\times\A_f=\R\times\A_f^p\times\Q_p^\times,
\qquad\A_f=\Z_f\otimes\Q,\qquad
\Z_f=\varprojlim_m\Z/m\Z=\Z_f^p\times\Z_p.
\]
\footnote{译注：上述分解最后一项原书印为 $\Q_p^\times$；由于 $\A$ 此处指阿代尔环，按环的分解应为 $\Q_p$。}
$K$ 为 $G(\Z_f)$ 的足够小的开子群。$\Delta$ 为有理素数的乘积，使 $p\nmid\Delta$ 时，
$p$ 在 $F$ 中非分歧（unramified），$B$ 在 $F$ 的全部 $p$ 上方素位处分裂，且
$K=K^pG(\Z_p)$，其中 $K^p=K\cap G(\A_f^p)$。

$S_\C={}_KS_\C$ 为 $G,K$ 与第 1 节定义的映射
$h:\C^\times\to G(\R)$ 所定义的 $\C$ 上 Shimura 簇；故
$S_\C(\C)=G(\Q)\backslash G(\A)/K_\infty K$，其中 $K_\infty$ 为
$h$ 在 $G(\R)$ 中的中心化子。
若 $V=V(\Z)$ 为 $\Z$ 模，则对任意环 $R$，记 $V(R)=V\otimes_\Z R$。

\origsection{1}{$S_\C$ 作为模概形（$S_\C$ as a moduli scheme）}
回忆，域 $k$ 上的 阿贝尔簇是底层簇完备且连通的 $k$ 上代数群；其群运算必交换，簇必射影。
一维 阿贝尔簇就是椭圆曲线（elliptic curve），可用形如
\[
Y^2Z=X^3+aXZ^2+bZ^3,\qquad a,b\in k,\quad4a^3+27b^2\ne0
\]
的方程描述。\footnote{\textbf{译者注（本次校订）：}这一短 Weierstrass 方程的通用写法需假设 $\operatorname{char}k\ne2,3$。一般特征下应使用含 $a_1XY$、$a_3Y$、$a_2X^2$ 等项的 Weierstrass 方程。下文转到 $\C$ 上时，自然满足这一条件。}直接用方程处理高维 阿贝尔簇往往不便；在 $\C$ 上，可以使用复向量空间中格（lattices）的经典描述。
设 $V$ 是 $\C^g$ 中的格，即某个实基生成的子群，故 $V\otimes_\Z\R\simeq\C^g$。
则 $\C^g/V$ 是紧复解析流形，加法使之成为交换 Lie 群。$g=1$ 时，对应于 $V$ 的
Weierstrass $\wp$ 函数及其导数，定义嵌入
\[
z\longmapsto(\wp(z),\wp'(z),1):\C/V\hookrightarrow\mathbf P_\C^2,
\]
使 $\C/V$ 为射影平面的代数子集。因此它自动有代数簇结构，是 阿贝尔簇。
$g>1$ 时不再如此，因为 $\C^g/V$ 上的函数可能太少，无法将它嵌入射影空间。
其任意亚纯函数都是 $\C^g$ 上 theta 函数（theta functions）的商，所以它为代数簇，当且仅当有足够多的
 theta 函数。

按定义，$V$ 的 theta 函数是 $\C^g$ 上的全纯函数 $\theta$，满足对 $v\in V$，
\[
\theta(z+v)=\theta(z)\exp(2\pi i(L(z,v)+J(v))),
\]
其中 $L(z,v)$ 关于 $z$ 复线性，$J(v)$ 只依赖于 $v$。可证明 $L$ 关于 $v$ 可加，因而延拓成
$L:\C^g\times\C^g\to\C$，第一变量复线性，第二变量实线性。置
$E(z,w)=L(z,w)-L(w,z)$。则
\begin{enumerate}
\item[(a)] $E$ 实值、实双线性且交替（alternating）；
\item[(b)] $E$ 在 $V\times V$ 上取整数值；
\item[(c)] $(z,w)\mapsto E(iz,w)$ 对称且非负。
\end{enumerate}
对称性等价于对所有 $z,w$ 有 $E(iz,iw)=E(z,w)$；非负性指对所有 $z$ 有
$E(iz,z)\ge0$。

满足这些条件的形式称为 $V$ 的 Riemann 形式。已知，存在足够多函数定义
$\C^g/V$ 的射影嵌入，当且仅当存在非退化 Riemann 形式（于是 $E(iz,z)$ 正定）。
若 $g=1$，可将 $E(z,w)$ 定义为边 $Ow,Oz$ 张成的有向平行四边形面积与格的基本平行四边形面积的比。
这个形式总存在，且在相差整数倍的意义下唯一，因此在一维情形中通常不再专门注明。$g>1$ 时，非退化
Riemann 形式通常不存在；存在时，也不在相差整数倍的意义下唯一。
但由于 $\C^g/V$ 紧，一个 Riemann 形式所定义的代数结构（不含射影嵌入）与所选形式无关。

因此，给定 $\C^g$ 中一个容许非退化 Riemann 形式的格，就得到 阿贝尔簇。
反过来，从 $g$ 维 阿贝尔簇 $A$ 可恢复一个 $g$ 维复向量空间 $W$ 及其中容许此种形式的格 $V$。
$W$ 有几种等价的描述：它可以看作 $A$ 的 Lie 代数 $\Lie(A)$、单位元处切空间 $t_A$，或拓扑流形
$A(\C)$ 的万有覆盖空间（universal covering space）。$V$ 可描述为指数映射
$\exp:\Lie(A)\to A(\C)$ 的核，或 $A(\C)$ 的基本群；该群交换，因而等于
$H_1(A,\Z)$。始终将同构 $W/V\simeq A(\C)$ 看成来自正合列
\[
0\longrightarrow H_1(A,\Z)\longrightarrow t_A\xrightarrow{\exp}A(\C)\longrightarrow0.
\]
由于 $H_1(A,\Z)$ 为 $t_A$ 中的格，
$H_1(A,\R)=H_1(A,\Z)\otimes\R\simeq t_A$。所以 $A$ 由 $H_1(A,\Z)$ 以及
$H_1(A,\R)$ 上的复向量空间结构确定。

实向量空间 $V(\R)$ 上的复结构定义同态
$h:\C^\times\to\Aut_\R(V(\R))$，$h(z):v\mapsto zv$，而此复结构也由 $h$ 确定。
故 阿贝尔簇 $A$ 由对 $(H_1(A,\Z),h)$ 唯一确定，其中 $h$ 由
$H_1(A,\R)=t_A$ 的复结构定义；反过来，每个容许非退化 Riemann 形式的对 $(V(\Z),h)$ 都来自
阿贝尔簇。

令 $V(\Z)=H_1(A,\Z)$。$A$ 的有限阶点对应于某个整数倍属于 $V(\Z)$ 的
$V(\R)$ 元素。更准确地说，有限阶点群可等同于
$V(\Q)/V(\Z)\subset V(\R)/V(\Z)$。对整数 $m>0$，$m$ 挠点群
$A_m(\C)=m^{-1}V(\Z)/V(\Z)\simeq V(\Z/m\Z)
\simeq(\Z/m\Z)^{2\dim A}$。
定义
\[
T_fA=\varprojlim_mA_m(\C)=V(\Z_f),\qquad
T_lA=\varprojlim_mA_{l^m}(\C)=V(\Z_l)
\]
（$l$ 为素数）。故 $T_fA=\prod_lT_lA$，$T_lA\simeq\Z_l^{2\dim A}$。

阿贝尔簇同态 $A\to A'$ 诱导复线性映射 $t_A\to t_{A'}$，将
$H_1(A,\Z)$ 映入 $H_1(A',\Z)$。反过来，若二者分别对应于
$(V,h)$、$(V',h')$，则 $\Z$ 模映射 $\alpha:V(\Z)\to V'(\Z)$，只要其延拓
$\alpha\otimes1:V(\R)\to V'(\R)$ 复线性（即与所有 $h(z),h'(z)$ 交换），就来自复流形映射
$A(\C)\to A'(\C)$；紧性又保证映射为代数的。用 $\End(A)$ 表示自同态环，
$\End^0(A)=\End(A)\otimes_\Z\Q$。它在 $H_1(A,\Q)$ 上有忠实表示，故为有限维
$\Q$ 代数；它还是半单代数，其可能的维数与结构已有完整的分类理论。

当 $A$ 对应于 $(V,h)$ 时，定义同态 $i:O_B\to\End(A)$，等价于在 $V$ 上定义
$O_B$ 作用，使 $h$ 将 $\C^\times$ 映入 $\Aut_{O_B\otimes\R}(V(\R))$。
给定这样的 $i$，且 $i(1)=1$ 时，就说 $O_B$ 作用于 $A$。它诱导单射
$i:B\hookrightarrow\End^0(A)$。

$A$ 的非退化 Riemann 形式 $E$，按规则
$E(\alpha z,w)=E(z,\alpha'w)$ 定义 $\End^0(A)$ 的对合
$\alpha\mapsto\alpha'$，称为 Rosati 对合。已知它正，即对 $\alpha\ne0$，
$\operatorname{Tr}(\alpha\alpha')>0$，这里 $\operatorname{Tr}$ 为从
$\End^0(A)$ 到 $\Q$ 的约化迹（reduced trace）。设 $O_B$ 作用于 $A$。
若存在非零 $c,c'\in O_F$，使对所有 $u,v\in V(\Z)$ 有
$E(u,cv)=E(u,c'v)$，就称两个 Riemann 形式 $E,E'$ 为 $F$ 等价（$F$-equivalent）；
并定义弱极化为非退化 Riemann 形式的一个 $F$ 等价类 $\Lambda$。
\footnote{译注：此处等价条件原书两边均印作 $E$；结合所定义的是 $E,E'$ 的等价关系，右边应涉及 $E'$。此处保留原印刷公式。}
由于 $F$ 是 $B$ 的中心，$\Lambda$ 中任意两个元素的 Rosati 对合在 $i(B)$ 上诱导相同映射。
我们考察满足
$E(i(b)u,v)=E(u,i(b^*)v)$（$u,v\in V(\R)$、$b\in B$、$E\in\Lambda$）的三元组
$(A,i,\Lambda)$；即要求 Rosati 对合稳定 $i(B)$，并诱导给定的 $b\mapsto b^*$。

回顾 $B$ 的一些记号。主对合（main involution）$b\mapsto b^\iota$ 定义为：在任一
$\R$ 同构 $B\otimes_F\R\simeq M_2(\R)$ 下，若
$b$ 对应 $M=\begin{pmatrix}a&b\\c&d\end{pmatrix}$，则 $b^\iota$ 对应
$M^\iota=\begin{pmatrix}d&-b\\-c&a\end{pmatrix}$。
于是 $b+b^\iota=\operatorname{Tr}_{B/F}(b)$，$bb^\iota=\operatorname{Nm}_{B/F}(b)$。
Skolem--Noether 定理保证存在 $t\in B$，使对所有 $b$，
$b^*=t^{-1}b^\iota t=tb^\iota t^{-1}$；自动有 $t^2\in F$，正性又推出
$t^2<0$，即其在每个 $F\hookrightarrow\R$ 下的像均负。
固定同构 $B\otimes_\Q\R\simeq M_2(\R)\times\cdots\times M_2(\R)$，使
$b\leftrightarrow(M_1,\ldots,M_d)$ 时 $b^*\leftrightarrow(M_1^{\mathrm{tr}},\ldots,M_d^{\mathrm{tr}})$。
由于
\[
M^{\mathrm{tr}}=\begin{pmatrix}0&-1\\1&0\end{pmatrix}^{-1}
M^\iota\begin{pmatrix}0&-1\\1&0\end{pmatrix},
\]
$t$ 对应于 $(c_1\begin{pmatrix}0&-1\\1&0\end{pmatrix},\ldots,
 c_d\begin{pmatrix}0&-1\\1&0\end{pmatrix})$，$c_i\in\R$，且可选 $t$ 使所有 $c_i>0$。

下面的引理说明，若 $O_B$ 作用于复流形 $\C^{2d}/V$，则有 Riemann 形式，
其 Rosati 对合在 $B$ 上诱导 $b\mapsto b^*$，且任意两个这样的形式 $F$ 等价；即与
$O_B$ 作用及给定对合相容的弱极化唯一。

\begin{originaltheorem}{引理 1.1（Lemma 1.1）}
设 $V=V(\Z)$ 为秩 $4d$ 的自由 $\Z$ 模，带有 $O_B$ 作用。则
$V(\Q)$ 上存在非退化交替形式 $\psi$，使：
\begin{enumerate}
\item[(a)] $u,v\in V(\Z)$ 时 $\psi(u,v)\in\Z$；
\item[(b)] 对所有非零 $u\in V(\R)$，$\psi(ut,u)<0$；
\item[(c)] 对 $b\in B$、$u,v\in V(\Q)$，$\psi(bu,v)=\psi(u,b^*v)$；
\item[(d)] 对 $V(\Q)$ 的任意 $B$ 自同构 $\alpha$，有 $\mu(\alpha)\in F^\times$，使
$\psi(\alpha u,\alpha v)=\psi(u,\mu(\alpha)v)$。
\end{enumerate}
此外，若 $\psi'$ 是另一个满足 (c) 的非退化交替形式，则存在
$c\in F^\times$，使 $\psi(u,cv)=\psi'(u,v)$。
\end{originaltheorem}
\begin{proof}
$V(\Q)$ 在 $B$ 上为一维，因此选合适基向量后，可将它等同于 $B$，将
$V(\Z)$ 等同于 $O_B$ 的左理想。定义
$\psi(u,v)=\operatorname{Tr}_{B/\Q}(uv^\iota t)=\operatorname{Tr}_{B/\Q}(utv^*)$。
则
\[
\psi(u,v)=\operatorname{Tr}(utv^*)=\operatorname{Tr}(vt^*u^*)
=\operatorname{Tr}(v(-t)u^*)=-\psi(v,u),
\]
所以 $\psi$ 是交替形式。(a) 显然\footnote{\textbf{译者注（本次校订）：}为得到 (a) 的整数值条件，可先将已选的 $t$ 乘以一个足够大的正整数，清除 $\psi$ 在格的有限组基向量对上的分母。这不改变对合，也不改变正性与非退化性。}，而非零 $u$ 时
\[
\psi(ut,u)=\operatorname{Tr}_{B/\Q}(ut^2u^*)
=\operatorname{Tr}_{F/\Q}(t^2\operatorname{Tr}_{B/F}(uu^*))<0,
\]
证明了 (b) 及非退化性。(c) 来自
$\psi(bu,v)=\operatorname{Tr}(utv^*b)=\operatorname{Tr}(ut(b^*v)^*)=\psi(u,b^*v)$。
最后，$B$ 的任意 $B$ 自同构均为右乘某个 $b\in B^\times$，故
$\psi(\alpha u,\alpha v)=\operatorname{Tr}(ubb^\iota v^\iota t)
=\psi(u,\mu(\alpha)v)$，其中 $\mu(\alpha)=\operatorname{Nm}_{B/F}(b)$。

为证末句，考察 $\Q$ 线性映射 $v\mapsto\psi'(1,v):B\to\Q$。
迹配对非退化，所以有唯一 $b\in B$，使对所有 $v$，
$\psi'(1,v)=\operatorname{Tr}(btv^*)$。于是
\[
\psi'(u,v)=\psi'(1,u^*v)=\operatorname{Tr}(btv^*u)
=\operatorname{Tr}(ubtv^*)=\operatorname{Tr}(ubv^\iota t).
\]
又
\[
\psi'(1,v)=-\psi'(v,1)=-\operatorname{Tr}(vbt)
=-\operatorname{Tr}(t^\iota b^\iota v^\iota)
=\operatorname{Tr}(b^\iota v^\iota t)=\operatorname{Tr}(b^\iota tv^*).
\]
故 $b=b^\iota$，于是 $b\in F$，可取 $c=b$。
\end{proof}

本节以后，取 $V(\Z)=O_B$，看成 $O_B$ 模，$\psi$ 同引理。
对任意环 $R$，可将 $G(R)$ 等同于 $\Aut_{O_B\otimes R}(V(R))$，因为
$O_B\otimes R$ 的每个自同态都是右乘其中元素。定义
$h:\C^\times\to G(\R)=\Aut_{B\otimes\R}(V(\R))$，使 $h(i)$ 为在
$B\otimes\R\simeq M_2(\R)^d$ 上右乘
$(\begin{pmatrix}0&1\\-1&0\end{pmatrix},\ldots,
\begin{pmatrix}0&1\\-1&0\end{pmatrix})$。于是 $K_\infty$ 由每个分量形如
$\begin{pmatrix}a&b\\-b&a\end{pmatrix}$ 的元组组成。
$E=\psi$ 是 $(V(\Z),h)$ 的 Riemann 形式，例如
\[
\psi(iu,iv)=\psi(uh(i),vh(i))
=\psi(u,\operatorname{Nm}_{B/F}(h(i))v)=\psi(u,v),
\]
且非零 $u$ 时 $\psi(iu,u)=\psi(ut/-(-t^2)^{1/2},u)>0$。
因此得到 阿贝尔簇 $A$，$O_B$ 作用诱导 $i:O_B\to\End(A)$，而包含 $\psi$ 的弱极化
$\Lambda$ 的 Rosati 对合在 $B$ 上诱导给定的 $*$。

回忆 $K\subseteq G(\Z_f)$ 为开子群。若存在 $k\in K$ 使
$\phi_1=k\phi_2$，就称同构 $\phi_1,\phi_2:T_fA\to V(\Z_f)$ 为
$K$ 等价（$K$-equivalent）。例如，若
$K=K_m=\ker(G(\Z_f)\to G(\Z/m\Z))$，给这样的同构的 $K$ 等价类，等价于给
$A_m(\C)\simeq V(\Z/m\Z)$，即一个 $m$ 级结构（level $m$ structure）。

\begin{originaltheorem}{定理 1.2（Theorem 1.2）}
点集 $S_\C(\C)=G(\Q)\backslash G(\A)/K_\infty K$ 与三元组
$(A,i,\overline\phi)$ 的同构类集一一对应，其中 $A$ 为 $2d$ 维 阿贝尔簇，
$i$ 定义 $O_B$ 在 $A$ 上的作用，$\overline\phi$ 为 $O_B$ 同构
$T_fA\simeq V(\Z_f)$ 的 $K$ 等价类。
\end{originaltheorem}

\paragraph*{注 1.3（Remark 1.3）。}
(a) 若存在同构 $\alpha:A\to A'$，使对所有 $b\in O_B$，
$\alpha\circ i(b)=i'(b)\circ\alpha$，且对每个
$\phi'\in\overline\phi'$ 有 $\phi'\circ T_f\alpha\in\overline\phi$，就称两个三元组同构。

(b) 通常研究 Shimura 簇参数化的 阿贝尔簇族，须使用四元组
$(A,i,\Lambda,\overline\phi)$，$\Lambda$ 为弱极化。本例无须如此，因为如已指出，
$\Lambda$ 总唯一存在。

\begin{proof}[定理 1.2 的证明]
先给每个 $g\in G(\A)$ 配三元组 $(A_g,i_g,\overline\phi_g)$。
$g=1$ 时取此前定义的 $(A,i)$，并取恒等映射
$T_fA=V(\Z_f)\xrightarrow{\mathrm{id}}V(\Z_f)$ 的等价类。
一般写 $g=g_\infty g_f$，分别用其无穷与有限部分改变 $V(\R)$ 的复结构和格。
定义 $h_g(z)=g_\infty h(z)g_\infty^{-1}$，并定义
$gV(\Z)=g_fV(\Z_f)\cap V(\Q)$，交在 $V(\A_f)$ 中进行。
取 $A_g$ 为对 $(gV,h_g)$ 定义的 阿贝尔簇。$O_B$ 仍作用于该格，故有显然的
$i_g:O_B\hookrightarrow\End(A_g)$。定义
$\phi_g:T_fA_g=g_fV(\Z_f)\to V(\Z_f)$ 为乘以 $g_f^{-1}$。

将 $g$ 换成 $gk_\infty$，$k_\infty\in K_\infty$，不改变 $h_g$；换成
$gk_f$，$k_f\in K$，也不改变复结构和格，只将 $\phi_g$ 换成与它 $K$ 等价的
$k_f^{-1}\phi_g$。换成 $qg$，$q\in G(\Q)$，则
$q^{-1}:V(\R)\to V(\R)$ 定义 $(A_{qg},\ldots)\simeq(A_g,\ldots)$。
故三元组只依赖于 $g$ 的双陪集。

反过来，同构 $\alpha:(A_g,\ldots)\to(A_{g'},\ldots)$ 由将
$gV(\Z)$ 同构映到 $g'V(\Z)$ 的 $V(\R)$ 同构诱导。特别地，它定义
$V(\Q)$ 的 $B$ 自同构 $q\in G(\Q)$。将 $g'$ 换成 $q^{-1}g'$、
$\alpha$ 换成 $q^{-1}\alpha$，即可假设其有理映射为恒等。
于是对所有 $z$，$g_\infty h(z)g_\infty^{-1}=g'_\infty h(z)(g'_\infty)^{-1}$，故
$g_\infty^{-1}g'_\infty\in K_\infty$。格相同又推出
$g_f^{-1}g'_f\in G(\Z_f)$；两个水平同构 $K$ 等价进一步推出
$g_f^{-1}g'_f\in K$。

最后证明这一对应是满射。由于 $B$ 为除代数，有
$B$ 同构 $H_1(A,\Q)\simeq V(\Q)$，据此将二者等同。$H_1(A,\Z)$ 为其中的格，故形如
$g_fV(\Z_f)\cap V(\Q)$，$g_f\in G(\A_f)$。同构 $V(\R)\simeq t_A$ 诱导复结构，记相应映射为
$h':\C^\times\to\Aut_\R(V(\R))$。由于 $B$ 在 $t_A$ 上复线性作用，$h'$ 落在
$G(\R)$ 中。显然存在 $g_\infty$ 使 $h'(z)=g_\infty h(z)g_\infty^{-1}$。
任意 $\phi\in\overline\phi$ 形如
$v\mapsto g_1^{-1}g_f^{-1}v:g_fV(\Z_f)\to V(\Z_f)$，$g_1\in G(\Z_f)$。
故取 $g=g_\infty g_fg_1$，即有 $(A,\ldots)\simeq(A_g,\ldots)$。
\end{proof}

\paragraph*{注 1.4（Remark 1.4）。}
(a) 阿贝尔簇映射 $\alpha:A\to A'$ 若满射且核有限，就称为同源（isogeny）；若与两个
$O_B$ 作用交换，就称为 $O_B$ 同源。复数上的同源显然诱导切空间同构，故
$A_g$ 与 $A_{g'}$ 同源\footnote{\textbf{译者注（本次校订）：}这里应指与给定 $O_B$ 作用相容的同源，即 $O_B$ 同源；下述 $G(\Q)$ 双陪集正是按 $B$ 线性有理同构分类。}，只有当 $g_\infty,g'_\infty$ 定义
$G(\Q)\backslash G(\R)/K_\infty$ 中同一双陪集时才可能。
另一方面，$\End_B^0(A)^\times\backslash G(\A_f)/K$ 分类存在 $O_B$ 同源
$A\to A_1$ 的三元组。例如 $g=g_f$ 时，乘以适当整数后可假设
$gV(\Z)\subset V(\Z)$；此时恒等映射 $V(\R)\to V(\R)$ 定义
$A_g\to A_1$ 的同源，核为 $V(\Z)/gV(\Z)$（参见第 6 节）。

(b) $F=\Q$ 时定理可加强。考虑投影
$V(\R)\times(G(\A)/K_\infty K)\to G(\A)/K_\infty K$。
给底空间通常的复结构，给 $gK_\infty K$ 上的 $V(\R)$ 副本以 $h_g$ 定义的复结构。
每个副本中有格 $gV(\Z)$，随 $g$ 连续变化。取商得到复流形映射
$\mathcal T\to G(\A)/K_\infty K$，其纤维为 $A_g$。
再令 $G(\Q)$ 作用于两边并取商，得到 阿贝尔簇解析族
$\mathcal A\to S_\C$，纤维带连续变化的 $(i_g,\overline\phi_g)$。
事实上这是代数族，即 $\mathcal A$ 为代数簇且映射代数。

若 $F\ne\Q$，上述构造失效，因为 $F$ 的单位可能作用于三元组，使群作用不自由。
但可作如下刚化（rigidification）：考虑四元组
$(A,i,\overline\phi,\varepsilon)$，前三项同前，$\varepsilon$ 为唯一弱极化
$\Lambda$ 到 $F^\times$ 的单射，满足
$\psi'(u,v)=\psi(u,cv)$ 时 $\varepsilon(\psi')=c\varepsilon(\psi)$。
这些同构类由 $F^\times\times S_\C(\C)$ 分类；它是每个 $F^\times$ 元素对应一个
$S_\C(\C)$ 副本的不交并，$F^\times$ 通过置换副本作用。
可将其看成 $\C$ 上无限个簇的不交并概形，前面的过程给出带结构的 阿贝尔簇代数族
$\mathcal A\to F^\times\times S_\C$。

\paragraph*{参考说明。}
关于复数上的 阿贝尔簇，初等介绍与更深入的讨论分别见 [20] 与 [17，第 I 章]。
Shimura 在 1960 年代深入研究了由 Shimura 簇参数化的 阿贝尔簇族（见其当时的
Annals 论文）。[4] 也有简短讨论。

\origsection{2}{$S$ 作为 $\Z[\Delta^{-1}]$ 上的概形（$S$ as a scheme over $\Z[\Delta^{-1}]$）}
下面将构造 $S_\C$ 的 $\Q$ 上模型 $S_\Q$，其基变换到 $\C$ 后同构于原簇。
这样的模型未必唯一；Shimura 给出了一组唯一刻画模型的条件，满足这些条件的模型称为典范模型。
例如，令 $F'$ 为使 $B$ 分裂的 $F$ 的全虚二次扩域，$A_0$ 为格
$O_{F'}\subset F'\otimes\R$ 定义的 $d$ 维 阿贝尔簇。$O_{F'}$ 作用于它，称它有来自
$F'$ 的复乘（complex multiplication）。令 $A=A_0\times A_0$。
将 $F'$ 嵌入 $B$，并选 $B$ 在 $F'$ 上的基 $\{e_1,e_2\}$，$e_i\in O_B$，得到
$B\hookrightarrow M_2(F')\subset M_2(\End^0(A_0))=\End^0(A)$，将
$O_B$ 送入 $\End(A)$。\footnote{\textbf{译者注（本次校订）：}有理代数的嵌入只直接给出 $B\hookrightarrow\End^0(A)$。仅要求 $F'$-基 $e_1,e_2$ 属于 $O_B$，还不足以保证 $O_B$ 保持 $O_{F'}e_1+O_{F'}e_2$，或保证下一式为格的同构。整性构造需选取受 $O_B$ 保持的格；必要时先将 $A_0^2$ 换成同源的阿贝尔簇，再配置水平结构。}还得到映射
\[
T_fA=(O_{F'}\oplus O_{F'})\otimes\Z_f\xrightarrow{\phi}
O_B\otimes\Z_f=V(\Z_f).
\]
三元组定义 $S_\C$ 的点，因而定义 $S_\Q$ 上一个复坐标点。模型典范时，这些坐标必须在
$\Q$ 上代数，并生成某个显式描述的类域（class field）。

为实现导言中的目标，需要构造由 $\Z[\Delta^{-1}]$ 中方程定义的概形 $S$，其有理纤维为典范模型，且每个
$p\nmid\Delta$ 时 $(S(\overline\F_p),\operatorname{Frob})$ 均可显式描述。
它定义函子 $R\mapsto S(R)$，给每个 $\Delta$ 可逆的环配以 $S$ 的
$R$ 坐标点集；更一般地，给每个 $\Spec\Z[\Delta^{-1}]$ 上的概形
$T$ 配以态射集 $S(T)=\Hom(T,S)$。此函子唯一确定概形，因此可以先写出函子
$\mathcal S$，使特别有 $\mathcal S(\C)=S_\C(\C)=G(\Q)\backslash G(\A)/K_\infty K$，再证明它为某个概形的点函子。

第 1 节提示，自然将 $\mathcal S(R)$ 定义为三元组 $(A,i,\overline\phi)$ 的同构类集，三项分别为相应复对象在
$R$ 上的类似物。$A$ 为 $R$ 上维数 $2d$ 的射影 阿贝尔概形（abelian scheme）。直观说，它为一族
阿贝尔簇，每个定义在 $R$ 的某个剩余域上；准确地说，它是 $\Spec R$ 上具有几何连通纤维的射影光滑群概形。
$i:O_B\hookrightarrow\End(A)$ 同前，满足 $i(1)=\mathrm{id}$，并假设
$A$ 有一个 Rosati 对合在 $B$ 上诱导 $*$ 的极化。

定义 $\overline\phi$ 时需处理两个问题。将 $\phi:T_fA\to V(\Z_f)$ 写为
$\prod_l\phi_l:\prod_lT_lA\to\prod_lV(\Z_l)$ 最易理解。
首先，若 $p$ 在 $R$ 中不可逆，就不可能有 $T_pA\simeq V(\Z_p)$，故只在
$R[p^{-1}]$ 上定义 $\phi_p$。其次，除非 $R$ 为代数闭域，不能期待每个 $l$ 都有这种同构，否则
$A$ 所有 $l$ 幂挠点的全部坐标都将属于 $R$。改为假设某个 $m$ 满足
$K\supset K_m=\ker(G(\Z_f)\to G(\Z/m\Z))$，并考察在 $R$ 的某个平展覆盖（étale covering）上定义的同构
$\phi:A_m\simeq V(\Z/m\Z)$。若在 $\Spec R$ 上局部有
$\phi_1=k\phi_2$，$k\in K$，就称二者 $K$ 等价，取此新意义下的等价类。
还须对三元组增加一个条件：若 $R$ 代数 $R'$ 满足
$O_B\otimes R'\simeq M_2(O_F\otimes R')$，则 $t_{A/R}$ 中对应于幂等元
$\begin{pmatrix}1&0\\0&0\end{pmatrix}$ 与 $\begin{pmatrix}0&0\\0&1\end{pmatrix}$ 的两个子模，在
$\Spec R'$ 上局部应为秩一自由 $O_F\otimes R'$ 模。$R=\C$ 时此条件自动成立；特征
$p$ 下类似情形中不成立的例子见 [18，1.29]。

接下来要判断这个函子是否可表。M. Artin 的可表性判据为此提供了可操作的方法：验证这些条件后，便可将函子实现为代数空间（algebraic space）的点函子。
代数空间比概形更一般，也足以处理这里的问题：它的环值点有明确意义，且平展上同调中容许通过模素位约化计算
Hasse--Weil $\zeta$ 函数的固有与光滑基变换定理（proper and smooth base change theorems）对它仍成立。
事实上，我们得到的代数空间几乎肯定是概形；使用 Mumford [16] 而非 Artin 的方法，大概就能证明。

先设 $F=\Q$。可核实 Artin 判据，得到 $\Z[\Delta^{-1}]$ 上固有光滑的代数空间
$S$，对每个 $\Delta$ 可逆的环，$S(R)=\mathcal S(R)$。特别地，
$S(\C)=\mathcal S(\C)=S_\C(\C)$，对 $p\nmid\Delta$ 有
$S(\overline\F_p)=\mathcal S(\overline\F_p)$。1.4(b) 的代数族
$\mathcal A\to S_\C$ 为 $\mathcal S(S_\C)=S(S_\C)$ 的元素，故给映射
$S_\C\to S$，诱导同构 $S_\C\to S\times\Spec\C$。而且已知 $S_\Q$ 是典范模型。

$F\ne\Q$ 时结果稍弱，但同样有用。三元组有非平凡自同构，所以不可能对所有
$R$ 都有 $S(R)=\mathcal S(R)$。不过存在同样固有光滑的代数空间，以及函子性映射
$\mathcal S(R)\to S(R)$，在 $R$ 为代数闭域时是同构。因此仍有复点集相等，且有理模型典范。
可按 1.4(b) 第二段刚化模问题，仿 $F=\Q$ 构造，再对 $F^\times$ 的左作用取商；或直接使用叠（stacks）。

两种情形中，$S(\overline\F_p)=\mathcal S(\overline\F_p)$ 都可用带附加结构的 阿贝尔簇描述。

\paragraph*{参考说明。}
[1] 简介 Artin 用代数空间表示函子的技术，[2] 更完整。[5]、[18] 将之用于与本题极相似的情形
（事实上几乎相同，见 [5] 导言第 7 节）。阿贝尔概形的基本定义见 [16]。

\origsection{3}{有限群概形、$p$ 可除群及 Dieudonné 模（Finite group schemes, $p$-divisible groups, and Dieudonné modules）}
以下须考虑 阿贝尔簇的有限子群概形，先回顾相关概念与性质。固定特征
$p\ne0$ 的完美域（perfect field）$k$。设 $R$ 为有限 $k$ 代数，即有限维
$k$ 向量空间，$N=\Spec R$。对 $k$ 代数 $R'$，
$N(R')=\Hom_{k\text{-alg}}(R,R')$。若给每个 $N(R')$ 交换群结构，使
$R'\to R''$ 诱导的映射都是群同态，则称 $N$ 连同这族结构为 $k$ 上有限群概形（finite group scheme）。
如同仿射代数群，给这些群结构等价于给余乘法（comultiplication）
$R\xrightarrow mR\otimes_kR$。

\paragraph*{例 3.1（Example 3.1）。}
(a) 任意有限交换群 $M$ 都可视为 $k$ 上代数群，从而为有限群概形。
令 $R$ 为每个 $M$ 的元素对应一个 $k$ 副本的乘积，$N=\Spec R$。作为集合，
$N=M$；群律诱导 $R$ 的余乘法，进而给全部 $N(R')$ 相容群结构。
若 $R'$ 除 $0,1$ 外无幂等元，则 $N(R')=M$。

(b) $\mu_{p^n}=\Spec k[T]/(T^{p^n}-1)$，其
$R'$ 点 $\{\zeta\in R':\zeta^{p^n}=1\}$ 在乘法下为群。
若 $R'$ 无幂零元（nilpotents），特别是整环，则点集仅为 $\{1\}$。

(c) $\alpha_p=\Spec k[T]/(T^p)$，
$\alpha_p(R')=\{a\in R':a^p=0\}$。由 $(a+b)^p=a^p+b^p$，加法给出群结构。
若 $R'$ 无幂零元，点集也只有一个元素。

(d) $\Z/p\Z=\Spec k[T]/(T^p-T)$。若 $R'$ 除 $0,1$ 外无幂等元（例如整环），则
$(\Z/p\Z)(R')=\F_p$ 为 $R'$ 的素子域，在加法下成群。
这是 (a) 的特例，因为
$k[T]/(T^p-T)=k[T]/(T(T-1)\cdots(T-(p-1)))\simeq k^p$。

$N=\Spec R$ 的秩或阶（rank or order）为 $\dim_kR$。
(a) 的阶为 $|M|$，而 $\mu_{p^n},\alpha_p,\Z/p\Z$ 的阶分别为
$p^n,p,p$。同态 $N_1=\Spec R_1\to N_2=\Spec R_2$ 是
$k$ 代数同态 $R_2\to R_1$，要求它诱导的所有点集映射都是交换群同态。

以下只考虑阶为 $p$ 幂的有限群概形。基本事实如下。

\paragraph*{事实 3.2（Facts 3.2）。}
(a) 它们构成阿贝尔范畴，所以可像模范畴中一样取核、商等。

(b) $k$ 代数闭时，单对象（simple objects）只有 $\mu_p,\alpha_p,\Z/p\Z$。
即每个对象都有组成列（composition series），其商是这三者之一；不同类型单对象之间没有非零同态。

(c) 此范畴自对偶，有反变函子 $N\mapsto\widehat N$，称为 Cartier 对偶（Cartier dual），给出自身的反变等价。
更准确地说，有配对 $N\times\widehat N\to\Gm=\GL_1$，使对任意
$k$ 代数 $R$，诱导同构
$N(R)\simeq\Hom_R(\widehat N,\Gm)$ 与
$\widehat N(R)\simeq\Hom_R(N,\Gm)$。
例如 $(\Z/p\Z)^\wedge=\mu_p$，配对为 $(n,\zeta)\mapsto\zeta^n$；
$\widehat{\alpha_p}=\alpha_p$，配对为
\[
(a,b)\longmapsto\exp(ab)=1+ab+\cdots+\frac{(ab)^{p-1}}{(p-1)!}.
\]

(d) $\Hom(\Z/p\Z,\Z/p\Z)=\Z/p\Z$，
$\Hom(\mu_p,\mu_p)=\Z/p\Z$，$\Hom(\alpha_p,\alpha_p)=k$。
第一项显然，第二项由 Cartier 对偶性得出。最后一项中对应 $c\in k$ 的映射，在代数上为
$T\mapsto cT:k[T]/(T^p)\to k[T]/(T^p)$，在点上为 $a\mapsto ca$。

(e) 若 $0\to N'\to N\to N''\to0$ 正合，则
$\operatorname{order}(N)=\operatorname{order}(N')\operatorname{order}(N'')$。

设 $A$ 为 $k$ 上 阿贝尔簇。每个
$A_{p^n}=\ker(p^n:A\to A)$ 都为阶 $(p^n)^{2\dim A}$ 的有限群概形，与
$p\ne\operatorname{char}k$ 时相同。系统
$A_p\hookrightarrow A_{p^2}\hookrightarrow\cdots$ 称为
$A$ 的 $p$ 可除群或 Barsotti--Tate 群 $A(p)$。
更一般地，高（height）为 $h$ 的 $p$ 可除群是系统
$N=(N_1\xrightarrow{i_1}N_2\xrightarrow{i_2}N_3\to\cdots)$，使
$N_n$ 阶为 $p^{nh}$，且 $i_{n-1}$ 将 $N_{n-1}$ 等同于
$p^{n-1}:N_n\to N_n$ 的核。
例如 $\Q_p/\Z_p=(\Z/p\Z\to\Z/p^2\Z\to\cdots)$ 与
$\mu_{p^\infty}=(\mu_p\to\mu_{p^2}\to\cdots)$ 高为一，$A(p)$ 高为
$2\dim A$。同态 $\phi:N\to N'$ 是与过渡映射交换的一族
$\phi_n:N_n\to N'_n$。

\paragraph*{练习 3.3（Exercise 3.3）。}
设 $k$ 代数闭。对任意 阿贝尔簇 $A$，存在映射
$\phi_n:A_{p^n}\to\widehat{A}_{p^n}$，使充分大的 $n$ 均满足
$\ker\phi_n=\ker\phi_{n+1}$。推出 $A$ 的 $p$ 挠点不多于 $p^{\dim A}$ 个，并且等号成立时
$A(p)=(\Q_p/\Z_p)^{\dim A}\times(\mu_{p^\infty})^{\dim A}$。
这样的 阿贝尔簇称为普通的（ordinary）。

令 $W=W_k$ 为 $k$ 上的 Witt 向量环（ring of Witt vectors）：它是特征零的完备离散赋值环，极大理想由
$p$ 生成，剩余域为 $k$。存在唯一自同构 $a\mapsto a^{(p)}$，在 $k$ 上诱导
$p$ 次幂映射。若 $k=\overline\F_p$，则 $W$ 为 $\Q_p$ 的极大非分歧扩张
$\Q_p^{\mathrm{un}}$ 的整数环的完备化，该自同构由通常的 Frobenius 诱导。
令 $W[F,V]$ 为非交换多项式环，满足对所有 $a\in W$，
\[
FV=p=VF,\qquad Fa=a^{(p)}F,\qquad aV=Va^{(p)}.
\]
有反变函子 $N\mapsto DN$，将每个 $p$ 幂阶有限群概形配以其 Dieudonné 模，一个作为
$W$ 模长度有限的 $W[F,V]$ 模；它给出范畴的反等价。
$DN$ 作为 $W$ 模的长度等于 $N$ 的阶。
\footnote{译注：原书此处将长度直接写成群概形的阶。按上述阶的定义及紧随其后的例子，应为
$\operatorname{order}(N)=p^{\operatorname{length}_W(DN)}$。译文保留原陈述并指出这一差异。}
因此，有限群概形的运算可转化为非交换环 $W[F,V]$ 上的模运算。例如\footnote{\textbf{译者注（本次校订）：}这里 $F=1$、$V=1$ 是选定基向量上的记号。由于 $F$ 为 Frobenius 半线性、$V$ 为逆 Frobenius 半线性，在 $k$ 上实际分别是 $a\mapsto a^p$ 与 $a\mapsto a^{1/p}$；一般不能将它们理解为底层加法群上的恒等映射。}
\[
D(\mu_p)=W/pW=k;\qquad F\text{ 作用为 }0,\quad V\text{ 作用为 }1.
\]
\[
D(\alpha_p)=k;\quad F=0,\ V=0;\qquad
D(\Z/p\Z)=k;\quad F=1,\ V=0.
\]
若 $N$ 是幺幂的且 $pN=0$，则 $DN=\Lie(\widehat N)$；这里的 Lie 代数具有零括号，
并作为 $p$-Lie 代数，$F$ 是其 $p$ 次幂运算，$V$ 为零。更一般地，若 $N$ 幺幂且为
$p^n$ 所杀，则 $DN=\Hom(N,W_n)$，其中 $W_1=\Ga$，$W_n$ 是看作代数群的长度为
$n$ 的 Witt 向量（Witt vectors）。存在典范非退化 $W$ 双线性配对
\[
\langle\ ,\ \rangle:DN\times D\widehat N\longrightarrow W\otimes\Q_p/\Z_p,
\]
满足
\[
\langle Fm,n\rangle=\langle m,Vn\rangle^{(p)},\qquad
\langle Vm,n\rangle^{(p)}=\langle m,Fn\rangle.
\]
函子 $D$ 可以延拓到 $p$ 可除群（$p$-divisible groups）。若 $N=(N_1\to N_2\to\cdots)$，
则 $DN=\varprojlim DN_n$ 是有限自由 $W$ 模，其秩等于 $N$ 的高度（height）。
两个这样的群 $N,N'$ 同源，当且仅当存在具有有限核的满射 $N\to N'$，也当且仅当存在余核
长度有限的单射 $DN'\to DN$。令
\[
W'=W[1/p]=W\otimes_{\Z_p}\Q_p,\qquad
W'[F,F^{-1}]=W'[F,V]=W[F,V]\otimes_{\Z_p}\Q_p.
\]
置 $D'N=DN\otimes_{\Z_p}\Q_p$，并将它看作 $W'[F]$ 模。于是 $N,N'$ 同源当且仅当
$D'N\simeq D'N'$。

以 $\mathcal M$ 表示所得 $W'[F]$ 模的范畴。若 $k$ 代数闭，则对每个有理数
$\lambda\in[0,1]$，恰有一个单对象（simple object）。写 $\lambda=s/r$，其中
$s,r$ 互素，则它是
\[
D^\lambda=W'[F]/(F^r-p^s),\qquad \dim_{W'}D^\lambda=r.
\]
其自同态环是 $\Q_p$ 上次数为 $r^2$ 的唯一除代数。\footnote{\textbf{译者注（本次校订）：}这里的 $r^2$ 是在 $\Q_p$ 上的维数，中心除代数的次数为 $r$；仅给出维数并不能唯一确定除代数，还须给出由斜率决定的 Brauer 不变量。例如分母同为 $3$ 的斜率 $1/3$ 与 $2/3$ 对应互为反代数的两个不同除代数。}每个对象都是有限直和
\[
(D^{\lambda_1})^{m_1}\oplus\cdots\oplus(D^{\lambda_t})^{m_t},
\]
其中 $\lambda_i$ 互异。称 $\lambda_i$ 为斜率（slopes），其重数为 $m_ir_i$。
有时将 $(D^{s/r})^m$ 写成 $D^{sm/rm}$，因而分子、分母不一定互素：$D^{s/r}$ 的
斜率为 $s/r$，重数及维数均为 $r$。

因此，在代数闭域上，$D'N$ 的斜率给出 $N$ 的同源分类。高度为 $1$ 的群同构于
$\mu_{p^\infty}$ 或 $\Q_p/\Z_p$，且
\[
D'(\mu_{p^\infty})=D^1,\qquad D'(\Q_p/\Z_p)=D^0.
\]
二维单对象只有 $D^{1/2}=D'(A(p))$，其中 $A$ 是超奇异椭圆曲线（supersingular
elliptic curve），见第 5 节。对一般完美域 $k$，有
$DN_{\bar k}\simeq DN\otimes_{W_k}W_{\bar k}$。若 $k=\F_q$，$q=p^a$，则
$F^a$ 是 $W$ 线性的；写
\[
\det(T-F^a\mid DN)=\prod_{i=1}^{\operatorname{ht}N}(T-\alpha_i),
\]
则 $D'N_{\bar k}$ 的斜率是 $\{\operatorname{ord}_q\alpha_i\}$，这里赋值归一化为
$\operatorname{ord}_q q=1$。

对阿贝尔簇 $A$，记 $DA=DA(p)$。将坐标取 $q$ 次幂的 Frobenius $\pi$ 在 $DA$ 上
诱导 $F^a$。因此，Mumford 第 19 节中的多项式 $P_A(T)$ 就是
$\det(T-F^a\mid DA)$，并给出上述斜率。还可形成投有限群概形
\[
T_fA=\varprojlim A_n,\qquad T_lA=\varprojlim A_{l^n}.
\]
若 $l\ne\operatorname{char}k$ 且 $k$ 代数闭，则 $T_lA$ 是秩为 $2\dim A$ 的自由
$\Z_l$ 模，而且 $T_fA=T_f^pA\times T_pA$。对象 $\mathcal M$ 中在 $F,V$ 下稳定的
格给出 $p$ 可除群的同构分类。这一理论的入门介绍见 [6]。

\origsection{4}{作为阿贝尔簇族的 $S(\overline{\F}_p)$（The family of abelian varieties）}
固定素数 $p\nmid\Delta$。由第 2 节，$S(\overline\F_p)$ 的点对应于三元组
$(A,i,\overline\phi)$：$A$ 是 $\overline\F_p$ 上维数为 $2d$ 的阿贝尔簇，$i$ 是
$\mathcal O_B$ 的作用，$\overline\phi$ 是同构
\[
\phi:T_f^pA\longrightarrow V(\Z_f^p),\qquad
T_f^pA=\varprojlim_{p\nmid n}A_n(\overline\F_p)
\]
的 $K^p$ 等价类。$\phi_p$ 仅在可逆化 $p$ 之后才有定义，而
$\overline\F_p[p^{-1}]$ 是零环。$t_A$ 作为 $\mathcal O_B\otimes\overline\F_p$ 模
须满足如下条件：\footnote{\textbf{译者注（本次校订）：}第一行应为 $\mathcal O_B\otimes\overline\F_p\simeq M_2(\mathcal O_F\otimes\overline\F_p)\simeq\prod_{j=1}^{d}M_2(\overline\F_p)$。仅写一个 $M_2(\overline\F_p)$ 会在 $d>1$ 时丢失其余因子；两个幂等元应在这个矩阵代数中理解。}
\begin{equation}
\begin{gathered}
\mathcal O_B\otimes\overline\F_p\simeq M_2(\overline\F_p)\text{ 中幂等元 }
\begin{pmatrix}1&0\\0&0\end{pmatrix},\quad
\begin{pmatrix}0&0\\0&1\end{pmatrix}\text{ 所对应的子空间}\\
\text{均为秩 }1\text{ 的自由 }\mathcal O_F\otimes\overline\F_p\text{ 模。}
\end{gathered}\tag{4.1}
\end{equation}
若 $A$ 的方程为 $\sum a_{(i)}T^{(i)}$，则以系数 $a_{(i)}^p$ 定义 $A^{(p)}$。
相对 Frobenius（relative Frobenius）$F_A:A\to A^{(p)}$ 将坐标取 $p$ 次幂，是次数为
$p^{2d}$ 的纯不可分同源。其核 $A_F$ 是阶为 $p^{2d}$、在域上只有一个点的群概形，
其合成列仅出现 $\alpha_p$ 和 $\mu_p$。作为群有 $D(A)=D(A^{(p)})$，但恒等映射
$DA\to DA^{(p)}$ 是 $(p)$ 线性的，即将 $am$ 变为 $a^{(p)}m$。复合
\[
DA\xrightarrow{\mathrm{id}}DA^{(p)}\xrightarrow{DF_A}DA
\]
是左乘 $F$。因为 $F_A$ 在 $t_A$ 上为零，故
\[
t_A\simeq t_{A_F}\simeq D\widehat{A_F}\simeq(DA_F)^*
 =\bigl(\operatorname{Coker}(DA\xrightarrow F DA)\bigr)^*.
\]
因而可以在 $DA/FDA$ 上检验 (4.1)。直观地说，点 $P$ 的坐标是 $A$ 的系数；
$\operatorname{Frob}P$ 对应于 $(A^{(p)},i^{(p)},\overline\phi^{(p)})$，其中的结构
选得使 $F_A$ 为三元组之间的映射。

再来看 Hecke 对应（Hecke correspondences）。设 $g\in G(\A_f)$，且开子群
$K'\subset G(\Z_f)$ 满足 $g^{-1}K'g\subset K$。右乘 $x\mapsto xg$ 给出双陪集空间
之间的映射，并给出代数映射
\[
\mathcal T(g):{}_{K'}S_{\C}\longrightarrow{}_KS_{\C}.
\]
若点 $P$ 对应于 $(A',i',\overline\phi')$，则其像对应于 $(A,i,\overline\phi)$，
当且仅当存在 $\mathcal O_B$ 同源 $\alpha:A\to A'$，使对某个正整数 $m$ 下图交换：
\[
\begin{tikzcd}
T_fA\arrow[r,"T_f\alpha"]\arrow[d,"\phi"']&T_fA'\arrow[d,"\phi'"]\\
V(\Z_f)\arrow[r,"mg"]&V(\Z_f).
\end{tikzcd}
\]
在特征 $p$ 中，仅有 $G(\A_f^p)$ 的作用；此时要求 $\alpha$ 的次数与 $p$ 互素，
图中使用除 $p$ 外的 Tate 模，且 $m$ 与 $p$ 互素。这两个映射均由
\[
\mathcal T(g):{}_{K'}S\times\Spec\Z_{(p)}\longrightarrow{}_KS\times\Spec\Z_{(p)},
\qquad \Z_{(p)}=\{m/n\in\Q:(n,p)=1\}
\]
基变换而来。因此，将特征零三元组特殊化到特征 $p$，与上述对应相容。

\origsection{5}{同源类（The isogeny classes）}
仍固定 $p\nmid\Delta$。考虑对 $(A,i)$，其中 $A$ 是 $\overline\F_p$ 上维数为 $2d$
的阿贝尔簇，$i:B\to\End^0(A)$ 且 $i(1)=1$。两个阿贝尔簇的同源关系记为
$A\sim A'$；对的 $B$ 同源关系也用此记号。以 $\mathcal I_p$ 表示所有 $B$ 同源类，
并以 $(A,i)\otimes\Q$ 表示相应的类。有自然映射
\[
S(\overline\F_p)\longrightarrow\mathcal I_p,\qquad
(A,i,\overline\phi)\longmapsto(A,i)\otimes\Q.
\]
它是满射（见本节末尾）。以下分两步进行：第 5 节描述 $\mathcal I_p$，第 6 节描述各个纤维。Frobenius 和 Hecke 对应保持这些纤维。

诱导 $B$ 上给定对合 $*$ 的唯一弱极化（weak polarization），对每个 $n$ 给出
$A_n\times A_n\to\Gm$ 的 $F$ 等价类（Mumford，第 23 节），从而给出非退化斜配对
\[
\psi_l:T_lA\times T_lA\longrightarrow T_l\Gm\simeq\Z_l\quad(l\ne p),
\qquad \psi_p:DA\times DA\longrightarrow W.
\]
后一个配对满足前述 $F,V$ 相容性，并且 $\psi_p(bm,n)=\psi_p(m,b^*n)$。
我们将利用有限域上的分类定理。回忆，阿贝尔簇为单的，意指它在基域上没有非零真
阿贝尔子簇；任意阿贝尔簇同源于单阿贝尔簇之积。以 $\pi_A$ 表示 Frobenius。

\begin{originaltheorem}{定理 5.1（Theorem 5.1）}
(a) 若 $A$ 在 $\F_q$ 上为单的，则 $E=\End^0A$ 是以 $\Q[\pi]$ 为中心的除代数。
$\pi$ 是代数整数，其在每个复嵌入下的绝对值为 $q^{1/2}$。其局部不变量为
\[
\operatorname{inv}_v(E)=
\begin{cases}
1/2,&v\text{ 为实位},\\
0,&v\mid l,\ l\ne p,\\
\displaystyle\frac{\operatorname{ord}_v\pi}{\operatorname{ord}_vq}
 [\Q[\pi]_v:\Q_p],&v\mid p.
\end{cases}
\]
此外
\[
2\dim A=[\Q[\pi]:\Q]\,[E:\Q[\pi]]^{1/2}.
\]
$e=[E:\Q[\pi]]^{1/2}$ 是这些不变量的最小公分母；若 $m(T)$ 是 $\pi$ 在 $\Q$ 上的
极小多项式，则 $P_A(T)=m(T)^e$。

(b) 两个单阿贝尔簇 $A,A'$ 同源，当且仅当存在将 $\pi_A$ 映到 $\pi_{A'}$ 的同构
$\Q[\pi_A]\to\Q[\pi_{A'}]$。

(c) 每个满足 (a) 中条件的 $\pi$ 都是某个 $\F_q$ 上单阿贝尔簇 $A_\pi$ 的 Frobenius。

(d) 对 $\F_q$ 上任意阿贝尔簇 $A,B$ 和任意素数 $l$（包括 $p$），典范映射
\[
\Hom(A,B)\otimes\Z_l\longrightarrow\Hom(A(l),B(l))
 =\Hom(T_lA,T_lB)
\]
是同构。若 $l\ne p$，右端将 $A(l),B(l)$ 看作 $\Gal(\overline\F_q/\F_q)$ 模。
\end{originaltheorem}
\begin{proof}
(a) 的第一部分是 Weil 的 Riemann 假设，(c) 来自 Honda，其余来自 Tate。
见 [17]、[21]、[7]、[22]、[23]。
\end{proof}
例如，取 $\pi=p^a$，$q=p^{2a}$，便得到椭圆曲线 $A_{p^a}$，其 $\End^0$ 是 $\Q$ 上
恰在 $p,\infty$ 分歧的四元数代数。这些曲线称为超奇异的。若 $\Q[\pi]$ 有实位，
则 $A$ 是超奇异椭圆曲线，或者在 $\F_{q^2}$ 上成为两个超奇异椭圆曲线之积的同源簇。
写
\[
p\mathcal O_F=\mathfrak p_1\cdots\mathfrak p_m,
\qquad d_i=[\mathcal O_F/\mathfrak p_i:\F_p],\qquad d=\sum_i d_i,
\]
其中这些素理想互异。

\begin{originaltheorem}{命题 5.2（Proposition 5.2）}
$B$ 在 $\End^0A$ 中的中心化子为下述两种之一：

(a) $F$ 上的四元数代数 $B'$，恰在无穷位、$B$ 不分裂的位及 $d_i$ 为奇数的
$\mathfrak p_i$ 处不分裂，而在其他位分裂；

(b) 使 $B$ 分裂的、$F$ 的全虚二次扩张 $F'$。

在 (a) 中，$A\sim A_0^{2d}$，其中 $A_0$ 是超奇异椭圆曲线；在 (b) 中，
$A\sim A_0^2$，且 $F'\subset\End^0A_0$。
\end{originaltheorem}
\begin{proof}
先设 $A\sim A_0^r\times A_1$，$r\ge1$，其中 $A_0$ 超奇异且 $\Hom(A_0,A_1)=0$。
于是 $\End^0A\simeq M_r(E)\times\End^0A_1$，$E=\End^0A_0$，且 $B$ 嵌入
$M_r(E)$。由 $F\hookrightarrow M_r(E)$，有 $d\mid2r$；但 $d=2r$ 不可能，
因为 $F$ 不使 $E$ 分裂（[19]，定理 10），故 $r=d$ 或 $2d$。\footnote{\textbf{译者注（本次校订）：}从 $d\mid2r$ 到这一结论还需补充：$E\otimes F$ 在实位仍为除代数，指标为 $2$；$F$ 在 $M_r(E)$ 中的中心化子的次数为 $2r/d$，且属于同一 Brauer 类，故 $2\mid2r/d$，即 $d\mid r$。结合 $1\le r\le2d$，才得到 $r=d$ 或 $2d$。}
由 Skolem--Noether 定理，可以将 $F$ 的映射共轭成经由 $M_r(\Q)$ 的映射；其中心化子为
\[
C(F)\simeq M_{r/d}(F)\otimes E=M_{r/d}(E\otimes F).
\]
若 $C$ 为 $B$ 在 $M_r(E)$ 中的中心化子，则 $B\otimes_FC\simeq C(F)$，且两者均为
中心单代数（central simple algebra）（[19]，第 8 节）。或者 $r/d=1$，此时 $C=F$、
$B=E\otimes F$；但 $B$ 在无穷位分裂而 $E\otimes F$ 不分裂，矛盾。或者 $r/d=2$，
此时 $C$ 是四元数代数，且在 Brauer 群中
$[B]+[C]=[E\otimes F]$，这就给出 (a)。

现在假定对超奇异 $A_0$ 有 $\Hom(A_0,A)=0$。选取足够大的有限域 $\F_q$，使 $A$
及其所有自同态均在其上有定义，并令 $\pi\in\End^0A$ 为 Frobenius。不存在
$\Q[\pi]\to\R$ 的同态。考虑各包含关系组成的图
\[
\begin{tikzcd}[column sep=large]
B\arrow[r,hook]&B[\pi]\arrow[r,hook]&B\otimes_FC\arrow[r,hook]&E\\
F\arrow[u,hook]\arrow[r,hook]&F[\pi]\arrow[u,hook]\arrow[r,hook]&C\arrow[u,hook]\\
\Q\arrow[u,hook]\arrow[r,hook]&\Q[\pi]\arrow[u,hook]
\end{tikzcd}
\]
其中 $E=\End^0A$，$C$ 为 $B$ 的中心化子。若 $F[\pi]=F$，便与假设矛盾。
又 $[F[\pi]:F]\le2$ 且 $F[\pi]=C$，否则 $E$ 将包含维数大于
$4d=2\dim A$ 的交换子环，与定理 5.1(a) 矛盾。
因此 $F'=C=F[\pi]$ 是没有实位的二次扩张。它使 $B$ 分裂：对有限素数 $l\ne p$，
$T_lA\otimes\Q_l$ 是 $F'_l=F'\otimes\Q_l$ 上秩为 $2$ 的自由模，故
$B\otimes F'_l\simeq M_2(F'_l)$；而 $B$ 本已在无穷位和 $p$ 处分裂。
取 $(B\otimes F')\cap\End A$ 中的幂等元 $e\ne0,1$，置 $A_0=eA$\footnote{\textbf{译者注（本次校订）：}这里可先在 $B\otimes F'\simeq M_2(F')\subset\End^0(A)$ 中取非平凡幂等元 $e$，再取正整数 $N$ 使 $Ne\in\End(A)$，并令 $A_0=\operatorname{Im}(Ne)$。一般不能预先要求有理幂等元本身属于 $\End(A)$；同源分解已足以完成论证。}，便有
$A\sim A_0\times A_0$。因 $F'$ 与 $e$ 交换，$F'\subset\End^0A_0$。
\end{proof}

\paragraph*{注 5.3（Remark 5.3）.}
在 (b) 中，$A_0\sim A_1^r$，其中 $A_1$ 为单阿贝尔簇，因为 $\End^0A_0$ 的中心是
$F'$ 的子域。因此任何这样的 $A$ 都同源于一个单阿贝尔簇的幂，$\End^0A$ 是
$\Q[\pi]$ 上的中心单代数。若 $\alpha:A\to A'$ 是同源，则 Skolem--Noether 定理
表明，$B\to\End^0A\xrightarrow{\alpha_*}\End^0A'$（其中
$\alpha_*(r)=\alpha r\alpha^{-1}$）与 $i'$ 相差一个内自同构 $\beta$，故 $\beta\alpha$
是 $B$ 同源。因而 $A\sim A'$ 蕴含 $(A,i)\sim(A',i')$。

进一步考察 (b)。设 $\mathfrak p_1,\ldots,\mathfrak p_t$ 在 $F'$ 中分裂，
$\mathfrak p_i=\mathfrak q_i\mathfrak q_i'$（$i\le t$），$0\le t\le m$。
$\mathcal O_F\cap\End A_0$ 与 $\Z_p$ 均作用于 $A_0(p)$，而
$F\otimes\Q_p=\prod F_{\mathfrak p_i}$ 诱导
\[
A_0(p)\sim\prod_iA_0(\mathfrak p_i),\qquad
D'A_0\simeq\prod_iD'A_0(\mathfrak p_i).
\]
每个因子的高度为 $2d_i$，相应 $D'$ 的维数也为 $2d_i$。由 $\psi_p$ 下 $F$ 的自伴性，
此分解正交，且每个因子上的配对非退化。因此其斜率在 $\lambda\mapsto1-\lambda$ 下不变。
当 $i\le t$ 时，$F'_{\mathfrak p_i}\simeq F'_{\mathfrak q_i}\times F'_{\mathfrak q_i'}$
的作用进一步给出
\[
A_0(\mathfrak p_i)\sim A_0(\mathfrak q_i)\times A_0(\mathfrak q_i'),\qquad
D'A_0(\mathfrak p_i)\simeq D'A_0(\mathfrak q_i)\times D'A_0(\mathfrak q_i').
\]
因为 $[F'_{\mathfrak q_i}:\Q_p]=d_i=\operatorname{ht}A_0(\mathfrak q_i)$，
$A_0(\mathfrak q_i)$ 同源于一个单 $p$ 可除群的幂，故
\[
D'A_0(\mathfrak q_i)=D^{k_i/d_i},\qquad
D'A_0(\mathfrak q_i')=D^{k_i'/d_i},\qquad
0\le k_i\le d_i,\quad k_i+k_i'=d_i.
\]
当 $i>t$ 时，$[F'_{\mathfrak q_i}:\Q_p]=2d_i=\operatorname{ht}A_0(\mathfrak p_i)$，
因而它仍同源于一个单群的幂。若写其 Dieudonné 模为 $D^{s/r}$，斜率的对称性强迫
$s/r=1/2=d_i/(2d_i)$；此时置 $k_i=d_i/2$。

至少有一个 $k_i\ne d_i/2$。否则所有斜率均为 $1/2$，在所有位都有
$|\pi/q^{1/2}|_v=1$，所以某个幂是单位根。扩大 $\F_q$ 后可令 $\pi=q^{1/2}$，
因而处于超奇异情形 (a)，矛盾。特别地 $t\ge1$。

\begin{originaltheorem}{定理 5.4（Theorem 5.4）}
$\mathcal I_p$ 中，对每个满足下列条件的对 $(F',(k_i)_{1\le i\le m})$ 恰有一个元素：
$F'$ 是使 $B$ 分裂的 $F$ 的全虚二次扩张，至少一个 $\mathfrak p_i$ 在其中分裂；
若 $\mathfrak p_i$ 分裂，$k_i$ 为 $0\le k_i\le d_i$ 的整数，否则 $k_i=d_i/2$；
至少对一个 $i$ 有 $k_i\ne d_i-k_i$。当 $\mathfrak p_i$ 分裂时，$k_i,k_i'$ 分别
对应于 $\mathfrak q_i,\mathfrak q_i'$；须将 $F$ 上共轭的对
$(F',(k_i))$ 与 $(\overline F',(\overline k_i))$ 视为相同。此外还恰有一个
“超奇异”元素。
\end{originaltheorem}
例如，若 $F=\Q$，则除超奇异同源类外，对每个使 $B$ 分裂且 $p$ 在其中分裂的虚二次
数域 $F'$ 恰有一类。若 $p$ 在 $F$ 中完全分裂，则除超奇异类外，对每个具有适当类型的
全虚二次扩张 $F'/F$，以及对每个在 $F'$ 中分裂的 $\mathfrak p_i$ 上方的一对素理想所作
的二选一选择，恰有一类；一组选择与将每一项都换成其共轭后所得的选择视为相同。
\begin{proof}[定理 5.4 的证明]
先构造对应于 $(F',(k_1,\ldots,k_m))$ 的同源类 $(A,i)\otimes\Q$。仍令
$\mathfrak p_i$（$1\le i\le t$）为整除 $p$ 且在 $F'$ 中分裂的素理想。考虑
$\mathcal O_{F'}$ 中的理想\footnote{\textbf{译者注（本次校订）：}当某个未分裂的 $\mathfrak p_i$ 在 $F'/F$ 中分歧时，下式指数还需乘以分歧指数。具体地，对 $i>t$ 应取 $\mathfrak q_i^{e_i f/2}$，其中 $e_i=e(\mathfrak q_i/\mathfrak p_i)$；这样才有 $\mathfrak a\overline{\mathfrak a}=(p^f)$。正文指数 $f/2$ 适用于 $e_i=1$ 的惰性情形。紧接着的 $(\pi\bar\pi)$ 理想计算中，相应的 $\mathfrak q_i$ 指数也应为 $e_i f h$。}
\[
\mathfrak a=\mathfrak q_1^{(f/d_1)k_1}(\mathfrak q_1')^{(f/d_1)k_1'}
 \cdots\mathfrak q_{t+1}^{(f/d_{t+1})k_{t+1}}\cdots\mathfrak q_m^{(f/d_m)k_m},
\qquad f=2d_1\cdots d_m.
\]
对某个 $h$，$\mathfrak a^h$ 为主理想，记为 $(\pi)$。以 $a\mapsto\bar a$ 表示
$F'/F$ 的非平凡自同构，则 $\pi\bar\pi\in F$，且
\[
(\pi\bar\pi)=
 \bigl(\mathfrak q_1^f(\mathfrak q_1')^f\cdots\mathfrak q_m^f\bigr)^h\cap\mathcal O_F
 =\mathfrak p_1^{fh}\cdots\mathfrak p_m^{fh}=(p^{fh}).
\]
因此 $\pi\bar\pi=up^{fh}$，$u$ 为 $\mathcal O_F$ 的单位。若 $u$ 在 $F$ 中是平方，
以 $\pi/u^{1/2}$ 代替 $\pi$，得到 $\pi\bar\pi=q$，其中 $q=p^{fh}$。否则以
$\pi^2/u$ 代替 $\pi$，得到同样的等式，但 $q=p^{2fh}$。注意，对某个 $i$ 有
$k_i\ne k_i'$ 蕴含 $\pi\notin F$，所以 $F'=F[\pi]$。任意嵌入 $F'\hookrightarrow\C$
将 $F$ 映入 $\R$，因此复共轭在 $F'$ 上诱导 $a\mapsto\bar a$。特别地，$\bar\pi$
是复数 $\pi$ 的共轭，故 $\pi\bar\pi=q$ 蕴含 $|\pi|=q^{1/2}$。

令 $A_\pi$ 为定理 5.1(c) 中对应于 $\pi$ 的阿贝尔簇，置 $E=\End^0(A_\pi)$。
对 $\Q[\pi]$ 的任意素位 $v$，
\[
\operatorname{inv}_v(E)=
\begin{cases}
0,&v\nmid p,\\
(k_i/d_i)[\Q[\pi]_v:\Q_p],&v\mid p,\ \mathfrak q_i\mid v,\\
(k_i'/d_i)[\Q[\pi]_v:\Q_p],&v\mid p,\ \mathfrak q_i'\mid v.
\end{cases}
\]
以 $e_v$ 表示最简分数 $\operatorname{inv}_v(E)$ 的分母，以 $e$ 表示所有 $e_v$ 的最小
公倍数。于是 $2\dim A_\pi=re$，其中 $r=[\Q[\pi]:\Q]$。显然，对任何
$\mathfrak q\mid v$、$v\mid p$，有 $e_v\mid[F'_{\mathfrak q}:\Q[\pi]_v]$。
由类域论，这蕴含 $F'$ 使 $E$ 分裂；又显然 $e$ 整除
$[F':\Q[\pi]]=2d/r$。因为
\[
[M_{2d/re}(E):\Q[\pi]]=(2d/re)^2e^2=[F':\Q[\pi]]^2,
\]
$F'$ 嵌入 $M_{2d/re}(E)$（[19]，定理 10）。置 $A_0=A_\pi^{2d/re}$。
由 5.1(a)，$\pi$ 在 $A_\pi$ 上的特征多项式为 $c_\pi(T)^e$，其中 $c_\pi(T)$ 是
$\pi\in\Q[\pi]$ 在 $\Q$ 上的极小多项式。因此 $P_{A_0}(T)=c_\pi(T)^{2d/r}$，
恰为 $\pi\in F'$ 在 $\Q$ 上的特征多项式。对应于分裂
$F'_p=F'_{\mathfrak q_1}\times F'_{\mathfrak q_1'}\times\cdots$，有
\[
A_0(p)\sim A_0(\mathfrak q_1)\times A_0(\mathfrak q_1')\times\cdots,
\qquad P_{A_0}(T)=P_1(T)P_1'(T)\cdots,
\]
其中 $P_i(T)$、$P_i'(T)$ 分别为 $\pi$ 在 $F'_{\mathfrak q_i}$、$F'_{\mathfrak q_i'}$
中的像 $\pi_i,\pi_i'$ 在 $\Q_p$ 上的特征多项式。因此由第 3 节，$A(\mathfrak q_i)$
的斜率为 $\operatorname{ord}_q(\pi_i)=(fh/d_i)k_i/fh=k_id_i$，而
$A(\mathfrak q_i')$ 的斜率为 $k_i'/d_i$。
\footnote{译注：扫描页此处等式末项印作 $k_id_i$，而前面的商式及本节的记号给出
$k_i/d_i$；此处保留可见印式。}
于是，视为 $\overline\F_p$ 上阿贝尔簇的 $A=A_0\times A_0$，连同
$B\hookrightarrow M_2(F')$ 诱导的 $i$，代表对应于 $(F',(k_1,\ldots,k_m))$ 的同源类。
若 $A_0$ 为超奇异椭圆曲线，则 $(A_0^{2d},i)$ 代表超奇异类。

显然，若 $(A,i)\sim(A',i')$，则两者或者都代表超奇异类，或者对应于同一个对
$(F',(k_1,\ldots,k_m))$。尚须证明，若两者均对应于此对，则它们同源。将 $A,A'$
看作在 $\overline\F_p$ 的某个有限子域上有定义，得到
$\pi=\pi_A\in F'$ 与 $\pi'=\pi_{A'}\in F'$。假设蕴含，对每个 $\mathfrak q\mid p$
有 $\operatorname{ord}_q(\pi)=\operatorname{ord}_q(\pi')$；由 5.1(a)，在 $F'$ 的
所有素位上有 $|\pi/\pi'|_v=1$。故 $\pi,\pi'$ 相差一个单位根。扩大有限域后，可以
使它们相等。于是 $A,A'$ 同源于同一个 $A_\pi$ 的幂，因此彼此同源，最后应用 5.3
即完成证明。
\end{proof}

要证明 $\mathcal I_p$ 中每个类均由 $S(\overline\F_p)$ 的一个元素代表，须用下列引理。
\begin{originaltheorem}{引理 5.5（Lemma 5.5）}
设 $T\subset T_fA$ 且 $T_fA/T$ 有限。则存在同源 $\alpha:A'\to A$，使
$T_f\alpha$ 将 $T_fA'$ 同构地映到 $T$。
\end{originaltheorem}
\begin{proof}
有限性意味着，对所有充分大的 $n$，$T/nT\to T_fA/nT_fA$ 的余核 $N$ 与 $n$ 无关。
因此存在映射 $\phi:A_n=T_fA/nT_fA\to N$。将 $A'$ 定义为
$a\mapsto(\phi(a),a):A_n\to N\times A$ 的余核，并以 $(b,a)\mapsto na$ 定义
$\alpha:A'\to A$；于是 $(T_f\alpha)(T_fA')=T$。
\end{proof}
设 $(A,i)$ 代表 $\mathcal I_p$ 中的一个类，并置
$\mathcal O'=B\cap\End(A)$，这是 $B$ 的一个序（order）。把 $(A,i)$ 看作在足够大的
有限域 $\F_q$ 上有定义，则
\[
\End(A)\otimes\Z_l\simeq\End_{\F_q}(T_lA),\qquad
\mathcal O'\otimes\Z_l=(B\otimes\Z_l)\cap\End_{\F_q}(T_lA).
\]
对几乎所有 $l$，$\mathcal O'\otimes\Z_l=\mathcal O_B\otimes\Z_l$，这时取
$T_l=T_lA$；对其余 $l$，可选取 $T_lA$ 中有限指数的 $\mathcal O_B$ 稳定子模 $T_l$，使
\[
\End_{\F_q}(T_l)\cap(B\otimes\Z_l)=\mathcal O_B\otimes\Z_l.
\]
注意 $D(T_p/pT_p)=M$ 是在 $W$ 上长度有限的 $W[F,V]$ 模。可以选择 $T_p$，使
$M/FM$ 满足 (4.1)。令 $A'$ 如引理中对应于 $T=\prod_lT_l$，则 $A'$ 连同显然的 $i$
及某个 $\overline\phi$ 属于 $S(\overline\F_p)$，并代表 $(A,i)\otimes\Q$。

\origsection{6}{一个同源类（An isogeny class）}
尚须描述 $S(\overline\F_p)$ 中这样的元素 $(A',i',\overline\phi')$ 的集合
$Z=Z(A,i,\phi_A)$：$(A',i')$ 与给定的 $(A,i)$ 同源。一个 $\mathcal O_B$ 同源
$\alpha:A'\to A$ 确定单射 $T_f\alpha:T_fA'\to T_fA$，其像 $\Lambda$ 满足：

(a) $\Lambda$ 在 $\mathcal O_B$ 下稳定。

(b) $T_fA/\Lambda$ 是有限群概形。更确切地，对充分大的 $n$，
\[
\operatorname{Coker}(\Lambda/n\Lambda\to T_fA/nT_fA)
 =\operatorname{Coker}(A_n'\to A_n)=\ker(A'\xrightarrow\alpha A).
\]
(c) $D(\Lambda/p\Lambda)/FD(\Lambda/p\Lambda)$ 满足 (4.1)，因为
$\Lambda/p\Lambda=A_p'$，从而\footnote{\textbf{译者注（本次校订）：}下式右端应为 $DA'/F(DA')$，因为 $\Lambda/p\Lambda=A_p'$。同源类内的积分 Dieudonné 模会变化，不能直接用原来的 $DA$ 代替。}
\[
D(\Lambda/p\Lambda)/FD(\Lambda/p\Lambda)=DA/F(DA).
\]
考虑满足 (a)、(b)、(c) 的所有 $T_fA$ 的子对象 $\Lambda$。任何这样的 $\Lambda$ 均可
写成 $\Lambda=\Lambda^p\times\Lambda_p$，其中 $\Lambda^p$ 是 $T_f^pA$ 中的
$\Z_f^p$ 格（取 $\Z_f^p$ 模的通常意义），而 $\Lambda_p\subset T_pA$。令 $Y$ 为
对 $(\Lambda,\overline\phi)$ 的集合，其中 $\Lambda$ 如上，$\overline\phi$ 为同构
$\Lambda^p\xrightarrow\sim V(\Z_f^p)$ 的 $K$ 等价类。因
$(\Lambda,\overline\phi)$ 由 $((\Lambda^p,\overline\phi),\Lambda_p)$ 确定，可写
$Y=Y^p\times Y_p$。由 5.5，每个 $(\Lambda,\overline\phi)\in Y$ 都对应于配有同源
$\alpha:A'\to A$ 的三元组 $(A',i',\overline\phi')\in S(\overline\F_p)$，因此有满射
$Y\to Z\subset S(\overline\F_p)$。

对正整数 $n$，置 $n(\Lambda,\overline\phi)=(n\Lambda,n\overline\phi)$。定义
$Y\otimes\Q$ 为对 $(y,n)$（$y\in Y$，$n\in\Z_{>0}$）的集合，并在
$n'y=ny'$ 时将 $(y,n)$ 与 $(y',n')$ 等同。于是写
\[
Y\otimes\Q=(Y^p\otimes\Q)\times(Y_p\otimes\Q).
\]
其中 $Y^p\otimes\Q$ 可等同于 $(T_fA)\otimes\Q$ 中的 $\mathcal O_B$ 稳定格
$\Lambda$，连同同构 $\Lambda\xrightarrow\sim T_fA$ 的一个 $K$ 等价类。
$H(\Q)=\End^0_{\mathcal O_B}(A)^\times$ 作用于 $Y\otimes\Q$：给定 $\alpha\in H(\Q)$，
选择正整数 $m$ 使 $m\alpha$ 为 $A$ 的同源，并定义
\[
\alpha(\Lambda,\overline\phi,n)
 =\bigl(T_f(m\alpha)\Lambda,\overline{\phi T_f(m\alpha)^{-1}},mn\bigr).
\]

\begin{originaltheorem}{引理 6.1（Lemma 6.1）}
上述映射 $Y\to Z$ 诱导双射 $H(\Q)\backslash(Y\otimes\Q)\xrightarrow\sim Z$。
\end{originaltheorem}
\begin{proof}
$(\Lambda,\overline\phi,n)$ 与 $(\Lambda',\overline\phi',n')$ 映到
$S(\overline\F_p)$ 的同一个元素，当且仅当存在 $\mathcal O_B$ 同源
\[
\begin{tikzcd}
A'\arrow[r,shift left=.7ex,"\alpha"]\arrow[r,shift right=.7ex,"\alpha'"']&A
\end{tikzcd}
\]
以及 $\mathcal O_B$ 同构 $\phi_0:T_fA'\to V(\Z_f)$，使\footnote{\textbf{译者注（本次校订）：}下面两式的水平结构应使用 $\phi_0\circ(T_f\alpha)^{-1}$ 与 $\phi_0\circ(T_f\alpha')^{-1}$，定义在相应的像格上；按原式直接复合，定义域与值域不相容。此处的水平结构还应统一使用除 $p$ 外的 Tate 模，正如本节开头的定义。}
\[
\begin{aligned}
n\bigl((T_f\alpha)T_fA',\overline{\phi_0(T_f\alpha)}\bigr)&=(\Lambda,\overline\phi),\\
n'\bigl((T_f\alpha')T_fA',\overline{\phi_0(T_f\alpha')}\bigr)&=(\Lambda',\overline\phi').
\end{aligned}
\]

于是 $\alpha'\alpha^{-1}$ 可看作 $\End^0(A)$ 的元素，且
$\alpha'\alpha^{-1}(\Lambda,\overline\phi,n)=(\Lambda',\overline\phi',n')$。
\end{proof}

\begin{originaltheorem}{引理 6.2（Lemma 6.2）}
映射
\[
G(\A_f^p)\longrightarrow Y^p\otimes\Q,\qquad
 g\longmapsto\bigl(g(T_f^pA),\overline\phi_Ag^{-1}\bigr)
\]
诱导双射 $G(\A_f^p)/K\to Y^p\otimes\Q$。\footnote{\textbf{译者注（本次校订）：}商中的子群应为 $K^p=K\cap G(\A_f^p)$，与定理 6.4 一致。}
\end{originaltheorem}
\begin{proof}显然。\end{proof}

\begin{originaltheorem}{引理 6.3（Lemma 6.3）}
$Y_p\otimes\Q$ 与如下集合 $X$ 一一对应：$D'A$ 中作为 $W$ 模自由、秩为 $4d$、
在 $\mathcal O_B$ 下稳定，且 $M/FM$ 满足 (4.1) 的 $W[F,V]$ 子模 $M$。
\end{originaltheorem}
\begin{proof}
$p:\Lambda\to\Lambda$ 诱导映射
$i_n:\Lambda/p^n\Lambda\hookrightarrow\Lambda/p^{n+1}\Lambda$，从而定义 $p$ 可除群
$\Lambda(p)=(\Lambda/p^n\Lambda,i_n)$。正合列
$0\to\Lambda\to T_pA\to N\to0$（$N$ 有限）给出
$0\to N\to\Lambda(p)\to A(p)\to0$。应用 $D$ 得
\[
0\longrightarrow DA\longrightarrow D\Lambda(p)\longrightarrow DN\longrightarrow0.
\]
因为 $DN$ 为挠模，可以将 $D'\Lambda(p)$ 与 $D'A$ 等同。对
$(\Lambda,n)\in Y^p$\footnote{译注：扫描页此处印为上标 $Y^p$；引理所述的是下标 $Y_p\otimes\Q$。\par\textbf{本次校订补充：}按此处的反变 Dieudonné 函子和前述分母记号，所配的格应为 $nD\Lambda(p)$。事实上，$(\Lambda,1)$ 与 $(m\Lambda,m)$ 表示同一元素，而在 $D'A$ 中有 $D(m\Lambda)(p)=m^{-1}D\Lambda(p)$；乘以分母后所得格才与代表元的选择无关。}，配以 $n^{-1}(D\Lambda(p))\in X$。

\end{proof}

\begin{originaltheorem}{定理 6.4（Theorem 6.4）}
沿用上述记号，有
\[
Z(A,i,\overline\phi)\simeq H(\Q)\backslash\bigl(G(\A_f^p)\times X\bigr)/K^p.
\]
$\operatorname{Frob}$ 将 $M\in X$ 送到 $FM$；Hecke 算子
$\mathcal T(g)$（$g\in G(\A_f^p)$）以 $G(\A_f^p)$ 上的右乘来“作用”。
\end{originaltheorem}
\begin{proof}汇总上述三个引理即得。\end{proof}
最后进一步描述 $X$。注意，对应于
$D'A\simeq D'A(\mathfrak p_1)\times\cdots\times D'A(\mathfrak p_m)$，有
$X\simeq X_1\times\cdots\times X_m$。方便起见，写
\[
\overline G_i(\Z_p)=\Aut_{\mathcal O_B}(A(\mathfrak p_i)),\qquad
\overline G_i(\Q_p)=\End_{\mathcal O_B}(A(\mathfrak p_i))^\times
 =\End_{\mathcal O_B}(D'A(\mathfrak p_i))^\times.
\]
在最简单的情形，$\overline G_i(\Q_p)$ 在属于 $X_i$ 的格
$M\subset D'A(\mathfrak p_i)$ 上传递作用，这时
$X_i\simeq\overline G_i(\Q_p)/\overline G_i(\Z_p)$。这里传递性意味着每个
$A'(\mathfrak p_i)$ 都同构于 $A(\mathfrak p_i)$，而不只是同源；参见 [6，第 93 页]。

\paragraph*{例 6.5（Examples 6.5）.}
(a) $F=\Q$，$F'$ 为 $\Q$ 的二次扩张，$p$ 在 $F'$ 中分解为
$(p)=\mathfrak q\mathfrak q'$，且 $(A,i,\overline\phi)$ 属于对应于 $(F',(0))$ 的同源类。
于是
\[
A(p)\simeq(\Q_p/\Z_p)^2\times(\mu_{p^\infty})^2,
\]
$\mathcal O_B\otimes\Z_p=M_2(\Z_p)$ 在各因子上按显然的方式作用。因此
\[
\overline G(\Z_p)=\left\{\begin{pmatrix}a&0\\0&b\end{pmatrix}:a,b\in\Z_p^\times\right\},
\qquad
\overline G(\Q_p)=\left\{\begin{pmatrix}a&0\\0&b\end{pmatrix}:a,b\in\Q_p^\times\right\}.
\]
此时 $X=\overline G(\Q_p)/\overline G(\Z_p)$，Frobenius 以
$\begin{pmatrix}1&0\\0&p\end{pmatrix}$ 作用。

(b) 条件同上，但 $(A,i,\overline\phi)$ 对应于 $(F',(1))$。于是
\[
A(p)\simeq(\mu_{p^\infty})^2\times(\Q_p/\Z_p)^2
\]
（即在 $F'_p=F'_{\mathfrak q}\times F'_{\mathfrak q'}$ 的分裂中，
$F'_{\mathfrak q}$ 现在对应于 $\mu_{p^\infty}$ 因子）。$\overline G(\Z_p)$、
$\overline G(\Q_p)$ 和 $X$ 与前相同，但 Frobenius 以
$\begin{pmatrix}p&0\\0&1\end{pmatrix}$ 作用。

(c) $F=\Q$，$(A,i,\overline\phi)$ 属于超奇异类。于是
$D'A(p)\simeq D^{1/2}\times D^{1/2}$，且 $\End(D^{1/2})=B_p'$ 是 $\Q_p$ 上唯一的
四元数除代数。$B$ 经由嵌入
\[
B\otimes\Q_p\simeq M_2(\Q_p)\overset{\mathrm{can}}{\hookrightarrow}M_2(B_p')
\]
作用。因此 $\overline G(\Q_p)$，即 $B\otimes\Q_p$ 在 $M_2(B_p')$ 中的中心化子，
等于 $(B_p')^\times$。而 $\overline G(\Z_p)$ 可取为 $\mathcal O^\times$，其中
$\mathcal O$ 是 $B_p'$ 的极大序（maximal order）。此时
$X\simeq\overline G(\Q_p)/\overline G(\Z_p)$。Frobenius 以乘上
$\widetilde\omega$ 作用，其中 $\widetilde\omega$ 是 $\mathcal O$ 的极大理想的生成元。

(d) $F$ 任意，$p$ 在 $F$ 中完全分裂，$(p)=\mathfrak p_1\cdots\mathfrak p_d$，
且 $(A,i,\overline\phi)$ 对应于 $(F',(k_1,\ldots,k_d))$。于是
$X\simeq X_1\times\cdots\times X_d$；若 $\mathfrak p_i$ 在 $F'$ 中分裂且 $k_i=0$，
则 $X_i$ 如 (a)；若分裂且 $k_i=1$，则如 (b)；否则如 (c)。

(e) 一般情形。结果的陈述见 [14]。这一情形在 J.~Milne，
\emph{Etude d'une classe d'isog\'enie}，S\'eminaire sur les groupes r\'eductifs et les
formes automorphes，Universit\'e Paris VII（1977--1978）中有详细讨论。

\paragraph*{校样增补（1978 年 11 月）.}
[12] 中对那些为模簇的 Shimura 簇的猜想所给出的证明纲要，没有会议当时看来那么完整。
上述对四元数代数乘法群情形的证明（在 6.5(e) 所引报告中补全）与该纲要略有区别，
它更倚重 Honda--Tate 关于有限域上阿贝尔簇同源类的分类。6.5(e) 所引的完整讨论班记录，
将以更详细的方式重述本文与 [3] 中的大量材料，并将在
\emph{Publications Math\'ematiques de l'Universit\'e Paris 7} 丛书中出版。

\begin{thebibliography}{23}
\bibitem[1]{c30-1} M. Artin, \emph{The implicit function theorem in algebraic geometry}, Colloq. in Algebraic Geometry, Bombay, 1969, pp. 13--34.
\bibitem[2]{c30-2} M. Artin, \emph{Th\'eor\`emes de repr\'esentabilit\'e pour les espaces alg\'ebriques}, Presses de l'Universit\'e de Montr\'eal, Montr\'eal, 1973.
\bibitem[3]{c30-3} W. Casselman, \emph{The Hasse--Weil $\zeta$-function of some moduli varieties of dimension greater than one}, 本论文集，第 2 册，pp. 141--163.
\bibitem[4]{c30-4} P. Deligne, \emph{Travaux de Shimura}, S\'eminaire Bourbaki, 1970/71, no. 389.
\bibitem[5]{c30-5} P. Deligne and M. Rapoport, \emph{Les sch\'emas de modules de courbes elliptiques}, Antwerp II, Lecture Notes in Math., vol. 349, Springer, New York, pp. 143--316.
\bibitem[6]{c30-6} M. Demazure, \emph{Lectures on $p$-divisible groups}, Lecture Notes in Math., vol. 302, Springer, New York, 1972.
\bibitem[7]{c30-7} T. Honda, \emph{Isogeny classes of abelian varieties over finite fields}, J. Math. Soc. Japan \textbf{20} (1968), 83--95.
\bibitem[8]{c30-8} Y. Ihara, \emph{The congruence monodromy problems}, J. Math. Soc. Japan \textbf{20} (1968), 107--121.
\bibitem[9]{c30-9} Y. Ihara, \emph{On congruence monodromy problems}, Lecture Notes, vols. 1, 2, Univ. Tokyo, 1968, 1969.
\bibitem[10]{c30-10} Y. Ihara, \emph{Non-abelian class fields over function fields in special cases}, Actes Congr. Internat. Math., vol. 1, Nice, 1970, pp. 381--389.
\bibitem[11]{c30-11} Y. Ihara, \emph{Some fundamental groups in the arithmetic of algebraic curves over finite fields}, Proc. Nat. Acad. Sci. U.S.A. \textbf{72} (1975), 3281--3284.
\bibitem[12]{c30-12} R. Langlands, \emph{Letter to Rapoport}（日期：1974 年 6 月 12 日至 9 月 2 日）.
\bibitem[13]{c30-13} R. Langlands, \emph{Some contemporary problems with origins in the Jugendtraum}, Proc. Sympos. Pure Math., vol. 28, Amer. Math. Soc., Providence, R.I., 1976, pp. 401--418.
\bibitem[14]{c30-14} R. Langlands, \emph{Shimura varieties and the Selberg trace formula}, Canad. J. Math. \textbf{29} (1977), 1292--1299.
\bibitem[15]{c30-15} Y. Morita, \emph{Ihara's conjectures and moduli space of abelian varieties}, 学位论文，Univ. Tokyo (1970).
\bibitem[16]{c30-16} D. Mumford, \emph{Geometric invariant theory}, Springer, 1965.
\bibitem[17]{c30-17} D. Mumford, \emph{Abelian varieties}, Oxford, 1970.
\bibitem[18]{c30-18} M. Rapoport, \emph{Compactifications de l'espace de modules de Hilbert--Blumenthal}, Compositio Math.（待刊）.
\bibitem[19]{c30-19} J. P. Serre, Expos\'e 7, S\'eminaire Cartan 1950/51.
\bibitem[20]{c30-20} H. Swinnerton-Dyer, \emph{Analytic theory of abelian varieties}, L.M.S. lecture notes \textbf{14} (1974).
\bibitem[21]{c30-21} J. Tate, \emph{Endomorphisms of abelian varieties over finite fields}, Invent. Math. \textbf{2} (1966), 134--144.
\bibitem[22]{c30-22} J. Tate, \emph{Classes d'isog\'enie des vari\'et\'es ab\'eliennes sur un corps fini}, S\'eminaire Bourbaki, 1968--1969, no. 352.
\bibitem[23]{c30-23} W. Waterhouse and J. Milne, \emph{Abelian varieties over finite fields}, Proc. Sympos. Pure Math., vol. 20, Amer. Math. Soc., Providence, R.I., 1971, pp. 53--64.
\end{thebibliography}
\begin{flushright}Michigan 大学（University of Michigan）\end{flushright}
